% =====================================================================
% SECTION III -- ENCORE: CERTIFIED ENVELOPE AND TWO-STAGE OFFERING
% =====================================================================
\section{Certified Envelopes and Two-Stage Offering}
\label{sec:methodology}

This section develops ENCORE for the problem of
Section~\ref{sec:problem_formulation}. The deliverable set is characterized
exactly as a context-dependent polyhedron whose maximal offer is the envelope
$F(x,c)$; nested conformal disturbance sets and exact tube margins robustify
it into $\Ft(x,c)$ (constructed in Section~\ref{subsec:tube}), from which the
day-ahead offer is drawn; and a certified
fallback policy underwrites the real-time layer, so the delivery guarantee
survives from commitment to metering. The remainder of the section
establishes exactness of the envelope, the virtual-battery equivalent, the
two-clause end-to-end guarantee, and the computational profile.

% --------------------------- FIGURE 2 --------------------------------
\begin{figure}[t]
\centering
\begin{tikzpicture}[font=\scriptsize, >=Stealth, node distance=3mm,
  box/.style={draw, rounded corners=1pt, align=center, inner sep=2.5pt},
  stage/.style={draw, rounded corners=2pt, inner sep=3pt, dashed}]
  % offline column
  \node[box, fill=yellow!12] (data) {trace + NWP archive\\ (fit block)};
  \node[box, fill=yellow!12, below=of data] (sets)
    {nested conformal sets\\ $\Wdel_{\!}(c)\subseteq\Wsafe_{\!}(c)$};
  \node[box, fill=yellow!12, below=of sets] (tube)
    {tube margins $M_t,\mu_t$\\ (gain $K$, ball $\bar e_0$)};
  \node[box, fill=yellow!12, below=of tube] (ftil)
    {tightened envelope $\Ft(x,c)$};
  \draw[->] (data) -- (sets);
  \draw[->] (sets) -- (tube);
  \draw[->] (tube) -- (ftil);
  % market column
  \node[box, fill=green!8, right=11mm of data] (offer)
    {day-ahead offer \eqref{eq:offer}\\ $q_h \le \Ft$, spaced};
  \node[box, fill=green!8, below=of offer] (act) {activation $r_h$};
  % online column
  \node[box, fill=blue!8, below=of act] (rt)
    {real-time: rolling LP\\ $\cup$ certified fallback\\
     $u_t=\bar u_t + K(x_t-\bar x_t)$};
  \node[box, fill=blue!8, below=of rt] (settle)
    {whole-hour settlement\\ \eqref{eq:settlement}};
  \draw[->] (ftil.east) -| ([xshift=-4mm]offer.south) -- (offer.south);
  \draw[->] (offer) -- (act);
  \draw[->] (act) -- (rt);
  \draw[->] (rt) -- (settle);
  \begin{scope}[on background layer]
  \node[stage, fit=(data)(sets)(tube)(ftil),
        label={[font=\scriptsize]above:offline (D$-$1)}] {};
  \node[stage, fit=(offer)(act)(rt)(settle),
        label={[font=\scriptsize]above:market + online (D-day)}] {};
  \end{scope}
\end{tikzpicture}
\caption{ENCORE pipeline. Offline, nested conformal sets and exact tube
margins tighten the deliverability polyhedron into $\Ft(x,c)$; the day-ahead
offer is drawn from its interior. Online, a certified fallback policy backs
the rolling controller, so Theorem~\ref{thm:main} survives to settlement.}
\label{fig:ctrl_arch}
\vspace{-3mm}
\end{figure}

% ---------------------------------------------------------------------
\subsection{The Deliverability Polyhedron and Its Envelope}
\label{subsec:envelope}

Fix a context $c$ and the product $(q,d)$ with $N=\dH/\Delta t$ control steps
per hour and an activation window of $m=d/\Delta t$ steps. Collecting the
discretized dynamics, the constraints C1--C4, and the two product clauses
into one system and eliminating the states gives a polyhedron in the lifted
variable $z=(x_0, q, u_{0:N-1})$:
\begin{equation}
\mathcal{D}(c) = \bigl\{ z : G z \le h(c) \bigr\},
\label{eq:lifted}
\end{equation}
with the context entering affinely through the floor \eqref{eq:floor}. Two
rows deserve emphasis because they encode the settlement of
Section~\ref{subsec:product}: the \emph{window-depth} rows
$\Pcool_t \le \Pbase - q$ for $t<m$ (the physical promise), and the
\emph{hour-energy} row
$\sum_{t<N}(\Pbase-\Pcool_t)\Delta t \ge r q \dH$ (the settled promise). The
pair is not redundant: with in-hour recovery, the window rows alone
over-promise, and the energy row alone would let depth be served by shallow
long reductions.

\begin{theorem}[Certified flexibility envelope]\label{thm:envelope}
For each context $c$: (i) $\mathcal{D}(c)$ is a polyhedron in $(x_0,q,u)$ and
the envelope $F(x,c)=\max\{q:\exists u,\ (x,q,u)\in\mathcal{D}(c)\}$ is
concave and piecewise affine in $x$; (ii) $F$ is nonincreasing in the
dew-point forecast and in the duration $d$; (iii) the projection of
$\mathcal{D}(c)$ onto any two coordinates of $(x_0,q)$ is computable exactly.
\end{theorem}

\begin{proof}
(i) Polyhedrality is by construction; $F$ is the optimal value of a linear
program whose feasible set is the polyhedral section of $\mathcal{D}(c)$ at
$x$, hence concave piecewise affine in the right-hand-side parameter $x$.
(ii) The forecast enters only through the floor \eqref{eq:floor}, which
tightens right-hand sides monotonically; raising $d$ adds window rows.
(iii) The two-dimensional shadow of a polyhedron is computed exactly by the
Bretl--Lall vertex-search procedure~\cite{bretl2008testing}.
\end{proof}

Define the \emph{ready state} $x^{\mathrm{ready}}(c)$ as the steady state at
nominal load with the supply temperature at the robust floor (the floor
\eqref{eq:floor} evaluated at the forecast shifted by the calibrated dew
margin of Section~\ref{subsec:sets}); it is the deepest pre-cool position the
condensation constraint certifiably allows, and the offering pre-positions
the plant there before each obligation. The startable sets
$\mathcal{R}(q)=\{x:F(x,c)\ge q\}$ are nested polytopes in the state plane;
Section~\ref{subsec:wall} computes them exactly on the calibrated hall.
Quantitatively, the calibrated model sustains $53$--$59\%$ of baseline
cooling power for $d=30$\,min from the nominal operating state, and $90\%$
from the pre-cooled ready state on a dry day; the envelope collapses toward
zero as the dew-point forecast approaches the supply floor. The thermal mass
is a genuine half-hour resource whose size is weather-coupled.

Projecting $\mathcal{D}(c)$ onto aggregate power--energy coordinates gives a
closed-form virtual-battery approximation~\cite{hao2015aggregate} whose
every parameter is a function of the context: state of charge
$\sigma = \sum_i C_i (T_i^{\lim} - T_i)/\eta$, capacity
\begin{equation}
E_{\mathrm{cap}}(c) = \frac{1}{\eta}\sum_i C_i\,\overline{\Delta T_i}(c),
\label{eq:vb}
\end{equation}
where $\overline{\Delta T_i}(c)$ is the admissible swing of node $i$ between
the ready state and its limit, a discharge limit set by the chiller share
net of the forced extraction $\ell(\Tw)$, a charge limit set by the
extraction headroom in \eqref{eq:Ux}, and leakage set by the loop exchange
rates. On the calibrated model, the coolant-loop tranche stores $45$\,kWh
per MW of IT load on a dry day (the facility loop adds the larger, slower
tranche that the three-state product certifies), every term shrinks affinely
as the dew-point forecast lifts the floor, and the closed forms track the
exact LP envelope within $3\%$ for durations $d\ge 15$\,min (median
$0.9\%$, worst case $2.7\%$). The object handed to an aggregator is
therefore \emph{a battery whose capacity depends on the weather}: a heat
wave that raises scarcity prices also raises the coastal dew point and
shrinks every hall's certified storage, which couples the resource's
availability to the price process and is captured by \eqref{eq:vb} at the
offering stage.

Finally, the envelope carries a structural property of energy-settled thermal
products.

\begin{proposition}[Commitment spacing]\label{prop:spacing}
Under whole-hour settlement, consecutive full-depth delivery is infeasible:
delivering $rq\dH$ in hour $h$ draws down the stored energy \eqref{eq:vb} by
at least $rq\dH$, so the terminal state lies outside $\mathcal{R}(q)$, and
restoring it during hour $h{+}1$ cancels that hour's settled energy.
\end{proposition}

ENCORE therefore separates committed hours by at least one uncommitted hour
as a design rule, sufficient for all depths, and uses the gap for
re-positioning (Lemma~\ref{lem:e0}).

% ---------------------------------------------------------------------
\subsection{Nested Context-Conditional Disturbance Sets}
\label{subsec:sets}

We now size the uncertainty the commitment must survive. Hourly disturbance
records are summarized by three physically chosen statistics: the peak step
residual, the positive residual energy, and the dew-point forecast residual.
For each context, $k$-nearest-neighbor split-conformal
calibration~\cite{lei2018distribution} of the three statistics yields a
\emph{budget polytope}
\begin{equation}
\mathcal{W}_{\varepsilon}(c)=\Bigl\{\xi:
\begin{aligned}
&\max_t |w_t| \le \bar w(c),\ \
|\delta^{\mathrm{dew}}|\le \bar D(c),\\[-1pt]
&\textstyle\sum_t [w_t]_+\Delta t \le \bar E(c)
\end{aligned}
\Bigr\},
\label{eq:box}
\end{equation}
with the miscoverage budget $\varepsilon$ allocated across the three faces a
priori (any allocation summing to $\varepsilon$ is valid by the union
bound). The amplitude face is calibrated on the upper tail, which dominates
this trace (negative residuals are verifiably smaller on the fit block, the
design split of Section~\ref{subsec:setup}); the
energy face is essential, because against bursty traces an amplitude-only
box is catastrophically conservative, while the budget face credits the fact
that a burst cannot persist. Coverage,
$\Pr\{\xi\in\mathcal{W}_\varepsilon(c)\}\ge 1-\varepsilon$, is
\emph{marginal} over hours exchangeable within the context neighborhood (the
deployed conditioning is hour-of-day; the dew-point forecast enters through
the floor and the ready state), and it is verified out of sample in
Section~\ref{sec:experimental_evaluation}. Commitment hours are selected by
value, and the validation therefore tests coverage on the \emph{selected}
hours directly rather than assuming it transfers.

Calibration uses $k=80$ conditioning neighbors and $150$ calibration points
per face, with hour-of-day features; the allocation $0.35/0.35/0.30$ across
the (amplitude, energy, dew) faces is fixed a priori, the dew share
reflecting the heavier-tailed NWP residuals. The two consequences of
uncertainty identified in Section~\ref{subsec:problem} have different
stakes, so they receive \emph{different} sets:
\begin{equation}
\Wsafe(c)=\mathcal{W}_{\epss}(c),\qquad
\Wdel(c)=\mathcal{W}_{\epsd}(c)\subseteq\Wsafe(c).
\label{eq:nested}
\end{equation}
Note the direction of the nesting: the \emph{smaller} miscoverage $\epss$
yields the \emph{wider} set, so $\Wdel\subseteq\Wsafe$ because
$\epsd>\epss$. The \emph{safety set} uses a small, fixed $\epss$ ($0.05$
here) and is never traded for revenue. The level $\epsd$ of the delivery set is a \emph{product
design variable} ($0.3$ at the design point): delivery failures beyond
$\Wdel$ are priced by the $\gamma$-penalty of \eqref{eq:settlement}, so
$\epsd$ trades certified depth against expected penalty and the offering
objective internalizes the trade. Beyond $\Wsafe$, a sub-$\epss$-probability
event, excursions are bounded by the plant's passive dynamics;
Section~\ref{sec:experimental_evaluation} reports the measured excursions.

% ---------------------------------------------------------------------
\subsection{Tube Tightening and the Fallback Certificate}
\label{subsec:tube}

Robustness enters the envelope through a fixed-gain error tube. Let $K$ be a
stabilizing gain (an offline linear-quadratic regulator design on the
discretized plant) and define
the \emph{fallback policy} $u_t=\bar u_t + K(x_t-\bar x_t)$ around a nominal
plan $(\bar x,\bar u)\in\mathcal{D}(c)$. The error $e_t=x_t-\bar x_t$ obeys
$e_{t+1}=(A+BK)e_t+Ew_t$ with $e_0$ in a pre-positioning ball, and the
componentwise worst-case margins over the budget polytope admit an
\emph{exact} greedy evaluation:
\begin{equation}
M_t=\!\!\max_{\substack{\xi\in\mathcal{W}_\varepsilon(c)\\ |e_0|\le\bar e_0}}\!\!|e_t|
 = \mathrm{greedy}\bigl(|(A{+}BK)^iE|;\ \bar w,\bar E\bigr)
 + |(A{+}BK)^t|\,\bar e_0,
\label{eq:margins}
\end{equation}
obtained by sorting the impulse-response coefficients and filling the energy
budget largest-first, plus the decayed initial-error term; the induced input
margins are $\mu_t=|K|M_t$. Each row of \eqref{eq:lifted} is then tightened
by the margins of the states and inputs it touches, \emph{with safety rows
(C1--C4) tightened by the $\Wsafe$-margins and the window-depth and
hour-energy rows by the $\Wdel$-margins}; the dew face shifts the floor,
$\vartheta(c)$ evaluated at $\Tdew^{\mathrm{fc}}+\bar D(c)$, in both the
constraint rows and the ready state. The result is the tightened envelope
$\Ft(x,c)\le F(x,c)$.

The ball $\bar e_0$ is not free: it is what the uncommitted hour of
Proposition~\ref{prop:spacing} can actually deliver.

\begin{lemma}[Disturbance-aware pre-positioning]\label{lem:e0}
Let the uncommitted hour preceding an obligation run the recovery law
$u=\bar Q + K_{\mathrm{rec}}(x-x^{\mathrm{ready}})$ elementwise on the loop
states, saturated into $\mathcal{U}(x;c)$, and suppose the realized
disturbance magnitude is bounded by $\bar w$ over the hour. With effective
gain $\kappa_{\mathrm{eff}}=\min\{K_{\mathrm{rec}},\,\mdot c_p\}$ (the
saturation binds near the floor) and step contraction $\rho<1$ away from
saturation, the start-of-hour error satisfies, componentwise,
\begin{equation}
|e_0| \;\le\; \rho^{N}\,|e_{\mathrm{init}}|
 \;+\; \bar w/\kappa_{\mathrm{eff}} .
\label{eq:e0}
\end{equation}
\end{lemma}

The binding term in \eqref{eq:e0} is the disturbance-driven floor, not the
transient: a $15$\,K post-event excursion contracts within the hour, while
the steady error never vanishes. The certificate's design ball
$\bar e_0 = 1.25$\,K instantiates \eqref{eq:e0} at the fit block's median
hourly disturbance magnitude with $K_{\mathrm{rec}}=80$\,kW/K and
$\kappa_{\mathrm{eff}}=\mdot c_p$; disturbances above that level can push a
start outside the ball, which is exactly what the validation of
Section~\ref{sec:experimental_evaluation} counts. A certificate built on
transient-only reasoning ($\bar e_0\approx 0.25$\,K) overstates readiness
and fails in closed loop.

\begin{theorem}[Two-clause end-to-end guarantee]\label{thm:main}
Let $q\le\Ft(x^{\mathrm{ready}},c)$ be committed for hour $h$ with the
spacing of Proposition~\ref{prop:spacing}, and let the hour run the fallback
policy from any $x_0$ with $|x_0-x^{\mathrm{ready}}|\le\bar e_0$. Then:
\emph{(i) Safety.} For every $\xi\in\Wsafe(c)$, constraints C1--C4 hold at
every step and the realized inputs remain admissible; hence safety holds
with probability at least $1-\epss$, marginally over context-hour draws.
\emph{(ii) Delivery.} For every $\xi\in\Wdel(c)$, the window depth and the
settled whole-hour energy obligations hold; hence the obligation is met with
probability at least $1-\epsd$, and, with the shortfall cap of
\eqref{eq:settlement}, the expected penalty is bounded by
$\epsd\,\gamma\,\picap_h r q \dH$.
\end{theorem}

\begin{proof}
By exactness of the support function of \eqref{eq:box} under the greedy
evaluation, $M_t$ in \eqref{eq:margins} upper-bounds $|e_t|$ for all
$\xi$ in the respective set and all $|e_0|\le\bar e_0$ (triangle inequality
for the initial-error term), and $\mu_t$ bounds the feedback's input
deviation. Feasibility of a tightened row therefore implies feasibility of
the original row along the realized trajectory: safety rows against
$\Wsafe$, delivery rows against $\Wdel\subseteq\Wsafe$, whose error bound is
also valid. The probability statements follow from the marginal validity of
the conformal sets \eqref{eq:nested}; the penalty bound prices the
complement of $\Wdel$ at the capped worst-case shortfall.
\end{proof}

The guarantee is per committed hour, conditional on the start premise
$|e_0|\le\bar e_0$; across a day, spacing decouples the hours and a union
bound compounds the miscoverage budgets over the committed hours.

\begin{remark}[Rolling improvement layer]\label{rem:mpc}
A rolling controller may improve on the fallback by re-solving the suffix of
\eqref{eq:lifted} from the measured state with margins restarted along the
tube, defaulting to the fallback on infeasibility or a missed solver
deadline. Inheriting Theorem~\ref{thm:main} additionally requires that the
improvement steps keep the state inside the certified tube around the
nominal plan; the evaluation of
Section~\ref{sec:experimental_evaluation} runs the certified fallback
itself during events, so the headline results carry the theorem's guarantee
directly.
\end{remark}

The guarantee never depends on the online layer: the certificate \emph{is}
the fallback. This inversion, a deliberately simple real-time layer
underwritten by an offline certificate, is the opposite of robustifying the
online controller itself as in~\cite{shen2026cdro}.

% ---------------------------------------------------------------------
\subsection{Day-Ahead Offering and Computational Profile}
\label{subsec:offering}

With hours decoupled by Proposition~\ref{prop:spacing} and the ready state
restored between obligations (Lemma~\ref{lem:e0}), the offer separates per
hour:
\begin{equation}
\begin{aligned}
\max_{0\le q\le \Ft(x^{\mathrm{ready}}_h, c_h)}\ \
&\picap_h q \dH - p_{\mathrm{act}}\bigl[r q \dH\,
\overline{\Delta\pi}^{\mathrm{rt}}_h + D_h(q)\bigr]\\
&\quad - p_{\mathrm{act}}\,\epsd\,\gamma\,\picap_h r q \dH,
\end{aligned}
\label{eq:offer}
\end{equation}
where the bracket carries the expected real-time cost of the energy-neutral
recovery shift and the degradation cost on the nominal plan, and the last
term prices the delivery risk that the certificate leaves to the settlement.
The objective is evaluated on a $q$-grid (exact up to grid resolution since
$D_h$ is convex in $q$), and spacing is enforced greedily by expected value.
Algorithm~\ref{alg:encore} summarizes the pipeline.

\begin{algorithm}[t]
\small
\caption{ENCORE Day-Ahead Certification and Offering}
\label{alg:encore}
\begin{algorithmic}[1]
\renewcommand{\algorithmicrequire}{\textbf{Input:}}
\renewcommand{\algorithmicensure}{\textbf{Output:}}
\REQUIRE context $c_h$ and prices for each hour; fit-block records; gain $K$;
levels $(\epss,\epsd)$; ball $\bar e_0$; product $(d, p_{\mathrm{act}},
\gamma)$
\ENSURE spaced commitments $\{q_h\}$ with certificates
\FOR{$h=1,\ldots,24$}
\STATE Calibrate $\Wsafe(c_h)\supseteq\Wdel(c_h)$ via \eqref{eq:box}--\eqref{eq:nested}
\STATE Compute margins \eqref{eq:margins} for both sets; tighten
\eqref{eq:lifted} $\to \Ft(\cdot,c_h)$
\STATE Pre-position target $x^{\mathrm{ready}}_h$ at the floor
\eqref{eq:floor} shifted by $\bar D(c_h)$
\STATE Solve \eqref{eq:offer} on a $q$-grid $\to (q_h, \text{value}_h)$
\ENDFOR
\STATE Enforce spacing (Proposition~\ref{prop:spacing}): greedy by value,
zero out adjacent commitments
\RETURN $\{q_h\}$; nominal plans and fallback gains as certificates
\end{algorithmic}
\end{algorithm}

The computational profile matches day-ahead practice. Every object in
Algorithm~\ref{alg:encore} is a small linear program in the lifted variable
($28$ decisions, $183$ rows at $N=12$): a full day's certification and
offering takes seconds on a workstation, and the online layer is a
fixed-gain affine policy plus an optional suffix LP re-solve, so no
iterative robust optimization runs in real time. The certified depth is
therefore limited by information (the context) rather than by computation.

\vspace{-2mm}
